\chapter*{پیشگفتار}
\phantomsection%
\addcontentsline{toc}{chapter}{پیشگفتار}%
\markboth{پیشگفتار}{پیشگفتار}
\label{sec:preface}

یادگیری ماشین به عنوان یکی از ارکان بنیادین هوش مصنوعی مدرن، در دهه‌های اخیر پیشرفت‌های مهمی را در نحوه تعامل سامانه‌های محاسباتی با داده‌ها رقم زده است. این شاخه از علم که بر آماره، بهینه‌سازی و یادگیری از نمونه‌ها استوار است، به سیستم‌ها امکان می‌دهد الگوهای پنهان را بدون طراحی صریح همه قواعد تصمیم‌گیری استخراج کنند \cite{726791,Vaswani2017}. با افزایش حجم داده‌ها و توان محاسباتی، یادگیری عمیق به چارچوب غالب بسیاری از کاربردهای هوش مصنوعی تبدیل شد و شبکه‌های عصبی عمیق توانستند در حوزه‌هایی مانند پردازش زبان طبیعی، تشخیص گفتار و بینایی ماشین عملکردی بسیار رقابتی ارائه دهند \cite{Vaswani2017,Dosovitskiy2020}. در بینایی ماشین، موفقیت شبکه‌های عصبی پیچشی در استخراج ویژگی از تصویر، نقطه شروع مهمی برای توسعه سامانه‌های دیداری مدرن بود و بعدها معماری‌های مبتنی بر مبدل نیز این مسیر را گسترش دادند \cite{726791,Dosovitskiy2020}.

بینایی ماشین، به عنوان یکی از مسیرهای اصلی کاربرد هوش مصنوعی در داده‌های بصری، از وظایف تحلیلی مانند طبقه‌بندی تصویر و تشخیص اشیا به سمت وظایف مولد نیز گسترش یافته است \cite{5206848,chen2025diffusionImageGenerationSurvey}. مدل‌های مولد نه تنها محتوای بصری را بازنمایی می‌کنند، بلکه می‌توانند نمونه‌های جدیدی مشابه توزیع داده‌های واقعی تولید کنند؛ این مسیر ابتدا با خانواده‌هایی مانند خودکدگذارهای تغییرپذیر\LTRfootnote{Variational Autoencoders} (\lr{VAE}) و شبکه‌های مولد تخاصمی\LTRfootnote{Generative Adversarial Networks} (\lr{GAN}) شکل گرفت و سپس با مدل‌های انتشار احتمالی به سطح تازه‌ای از کیفیت و پایداری رسید \cite{Kingma2013,Goodfellow2014,0214306}. مدل‌های انتشار با فرمول‌بندی تولید تصویر به عنوان وارون‌سازی تدریجی فرآیند نویزدهی، در بسیاری از وظایف تولید تصویر کیفیت و تنوع بالایی نشان داده‌اند و در مقایسه با \lr{GAN}ها از پایداری آموزشی مناسب‌تری برخوردارند \cite{8567753,chen2025diffusionImageGenerationSurvey}.

ترکیب مدل‌های انتشار با معماری‌های مبدل، به‌ویژه مبدل بینایی\LTRfootnote{Vision Transformer} (\lr{ViT})، مسیر شکل‌گیری مدل‌های مولد مقیاس‌پذیر را تقویت کرده است \cite{Dosovitskiy2020,2212.09748}. معماری‌هایی مانند \lr{DiT} با استفاده از بلوک‌های مبدل و سازوکارهایی مانند \lr{adaLN-Zero} برای شرط‌گذاری زمانی و کلاسی، نشان داده‌اند که افزایش ظرفیت و محاسبات می‌تواند کیفیت تولید تصویر را بهبود دهد \cite{2212.09748}. با این حال، استفاده از این مدل‌های بزرگ در دامنه‌های خاص، تنظیم دقیق کامل همه پارامترها (\lr{Full Fine-Tuning})، حافظه گرافیکی زیاد، زمان آموزش طولانی و ذخیره نسخه‌های متعدد از مدل را به مسئله‌ای پرهزینه تبدیل می‌کند \cite{5294562,8409325}.

در پاسخ به این چالش، روش‌های تنظیم دقیق پارامتر-کارا\LTRfootnote{Parameter-Efficient Fine-Tuning} توسعه یافته‌اند تا تنها بخش کوچکی از پارامترها برای سازگار کردن مدل با وظیفه جدید آموزش داده شود \cite{2479204,8409325}.
. با این حال، اگر تغییرات فقط در همان وضوح توکنی فشرده بدنه اصلی اعمال شوند، ممکن است بازسازی جزئیات مکانی و مؤلفه‌های فرکانس بالا محدود بماند. این مسئله انگیزه‌ای برای طراحی شاخه‌هایی است که بتوانند با وضوح مکانی بالاتر ویژگی‌های محلی را استخراج کنند و در عین حال بدنه اصلی مدل را ثابت نگه دارند \cite{wang2020hrnet,chen2021crossvit}.

یکی از راهکارهای مناسب برای این هدف، تفکیک وظایف «درک معنایی» و «حفظ جزئیات» در قالب معماری دوجریانی (\lr{Dual-Stream Architecture}) است \cite{wang2020hrnet,chen2021crossvit}. در این دیدگاه، جریان اصلی می‌تواند ساختار کلی و مفاهیم سطح بالا را حفظ کند، در حالی که جریان موازی کم‌هزینه روی جزئیات ریزدانه و ویژگی‌های فرکانس بالا تمرکز دارد \cite{wang2020hrnet}. چالش اصلی چنین طراحی‌ای، همگام‌سازی و ادغام کنترل‌شده اطلاعات جریان فرعی با بدنه اصلی است تا فرآیند تولید تصویر از نظر معنایی و مکانی منسجم باقی بماند \cite{chen2021crossvit,2209.12152}.

ترکیب هدفمند دو جریان نیازمند راهبردهایی فراتر از جمع ساده ویژگی‌هاست، زیرا خروجی شاخه وضوح‌بالا باید پیش از ادغام با جریان اصلی از نظر طول توالی و فضای ویژگی هم‌راستا شود \cite{chen2021crossvit,2209.12152}. استفاده از توجه متقاطع یا توکن‌های خلاصه‌ساز، از جمله توکن کلاس محلی، می‌تواند نقش یک واسط فشرده میان اطلاعات محلی و بازنمایی سراسری را ایفا کند \cite{Dosovitskiy2020,chen2021crossvit}. چنین سازوکاری به مدل اجازه می‌دهد در عین حفظ دانش کلی بدنه اصلی، اطلاعات ریزدانه و مکانی را نیز به‌صورت کنترل‌شده وارد فرآیند تولید کند \cite{2212.09748,2209.12152}.

علاوه بر ارتقای کیفیت بصری، منجمد نگه داشتن بخش اعظم شبکه و آموزش تنها شاخه کمکی می‌تواند فشار حافظه، تعداد پارامترهای آموزش‌پذیر و هزینه ذخیره مدل را کاهش دهد \cite{2479204,8409325}. این ویژگی برای کاربردهایی که نیازمند تنظیم دقیق سریع، شخصی‌سازی یا سازگاری با دامنه جدید هستند اهمیت ویژه‌ای دارد، زیرا امکان بهره‌گیری از \lr{Large Diffusion Models} را بدون آموزش مجدد کامل فراهم می‌کند \cite{ryu2023loradiffusion,0052338,2755774}. بنابراین، طراحی یک مسیر کمکی سبک در کنار بدنه اصلی منجمد، از نظر محاسباتی و کاربردی با جهت‌گیری پژوهش‌های اخیر در مدل‌های مولد همسو است \cite{2212.09748,8409325}.

در نهایت، تمرکز اصلی این پایان‌نامه بر اعمال تغییرات معماری هدفمند در مدل \lr{DiT} است تا تنظیم دقیق پارامتر-کارا با حفظ جزئیات محلی ممکن شود \cite{2212.09748,8409325}. بر اساس معماری نهایی پیاده‌سازی‌شده، یک شاخه موازی تحت عنوان $x2$ به شبکه افزوده می‌شود که ویژگی‌های مکمل را با وضوح مکانی بالاتر استخراج می‌کند و سپس از طریق توکن‌های کلاس محلی با جریان اصلی هم‌راستا و ادغام می‌شود \cite{chen2021crossvit,2088740}. ارزیابی این پژوهش روی تولید شرطی \lr{ImageNet-1K} نشان می‌دهد که \lr{DuoDiT-200} با آموزش \lr{17.95M} پارامتر از مجموع \lr{690.39M} پارامتر، به \lr{FID=2.219}، \lr{KID=0.0014} و \lr{LPIPS-Mean=0.7162} می‌رسد. این طراحی با منجمد نگه داشتن بدنه اصلی (\lr{Frozen Backbone})، امکان بهره‌گیری از دانش پیش‌آموزش‌دیده را حفظ می‌کند و هم‌زمان مسیر سبک‌تری برای سازگاری با داده‌ها یا کلاس‌های هدف فراهم می‌آورد \cite{2212.09748,Pan2024FineDiffusion}.
